% !TeX program = xelatex
% !TeX encoding = UTF-8
\documentclass[UTF8,a4paper,12pt]{ctexart}

\usepackage{geometry}
\usepackage{amsmath,amssymb,bm}
\usepackage{graphicx}
\usepackage{booktabs}
\usepackage{array}
\usepackage{longtable}
\usepackage{float}
\usepackage{enumitem}
\usepackage{hyperref}

\geometry{left=2.5cm,right=2.5cm,top=2.5cm,bottom=2.5cm}
\graphicspath{{../figures/}{figures/}}
\setlist{nosep}

\hypersetup{
  colorlinks=true,
  linkcolor=black,
  citecolor=black,
  urlcolor=blue
}

\title{\textbf{数学建模论文题目}}
\author{队伍编号：待填写}
\date{\today}

\begin{document}

\maketitle

\begin{abstract}
这里填写摘要。摘要应在完成建模、求解、验证和结果分析后再撰写，说明研究问题、模型方法、主要结果和结论。当前项目尚未提供数据和实验结果，因此不写入任何定量结论。
\end{abstract}

\noindent\textbf{关键词：}关键词一；关键词二；关键词三


\tableofcontents
\newpage

\section{问题重述与分析}

\subsection{问题重述}
本节用自己的语言说明赛题背景、研究对象、已知条件和待解决问题。正式写作时应避免照抄题面，并明确输入数据、约束条件和输出目标。

\subsection{问题分析}
本节分析问题的核心变量、约束关系和可选建模路线。若问题可以拆分为多个子问题，应说明子问题之间的依赖关系。

\begin{enumerate}
  \item 问题一：待根据赛题填写。
  \item 问题二：待根据赛题填写。
  \item 问题三：待根据赛题填写。
\end{enumerate}

\section{模型假设}

模型假设必须服务于赛题背景和数据条件，并在结果分析中说明其影响。当前没有具体赛题，因此仅保留占位。

\begin{enumerate}
  \item 假设一：待根据赛题背景填写。
  \item 假设二：待根据数据质量填写。
  \item 假设三：待根据模型适用条件填写。
\end{enumerate}

\section{符号说明}

符号表只列出正文中实际使用的符号。下表为格式占位，正式写作时请替换为真实定义。

\begin{table}[H]
  \centering
  \caption{符号说明占位表}
  \begin{tabular}{cl}
    \toprule
    符号 & 含义 \\
    \midrule
    $x_i$ & 第 $i$ 个样本或对象的输入变量，待根据赛题定义 \\
    $y_i$ & 第 $i$ 个样本或对象的目标变量，待根据赛题定义 \\
    $\bm{\theta}$ & 模型参数向量，待根据模型定义 \\
    $J(\bm{\theta})$ & 目标函数或损失函数，待根据模型定义 \\
    \bottomrule
  \end{tabular}
\end{table}

\section{数据预处理}

原始数据放入 \texttt{data/raw/}，不得直接修改。清洗后数据输出到 \texttt{data/processed/}，中间数据输出到 \texttt{data/interim/}。当前没有数据，因此不填写统计结果。

若需要对指标进行极差标准化，可使用
\begin{equation}
  \tilde{x}_{ij}
  =
  \frac{x_{ij} - \min_i x_{ij}}{\max_i x_{ij} - \min_i x_{ij}},
  \label{eq:min-max-normalization}
\end{equation}
其中 $x_{ij}$ 表示第 $i$ 个样本在第 $j$ 个指标上的原始取值。是否使用该式应由具体数据分布和建模目标决定。

\section{模型建立}

本节给出变量定义、目标函数、约束条件和建模理由。若问题可抽象为优化问题，可先写为通用形式：
\begin{equation}
  \min_{\bm{\theta} \in \Theta} J(\bm{\theta}),
  \label{eq:general-optimization}
\end{equation}
其中 $\Theta$ 为可行参数空间，$J(\bm{\theta})$ 为目标函数。正式论文中必须替换为与赛题一致的具体模型。

\section{模型求解}

本节说明求解算法、参数设置、计算环境和输出文件。所有可复现代码应放在 \texttt{code/} 中，所有图表应由代码输出到 \texttt{figures/} 或 \texttt{tables/}。当前没有数据和代码结果，因此不填写求解结论。

\section{模型检验与结果分析}

本节用于说明模型检验方法、误差指标、对照实验和结果解释。若使用预测模型，可根据具体任务定义误差指标，例如
\begin{equation}
  \mathrm{RMSE}
  =
  \sqrt{\frac{1}{n}\sum_{i=1}^{n}(y_i-\hat{y}_i)^2}.
  \label{eq:rmse}
\end{equation}
该指标只有在存在真实观测值 $y_i$ 和预测值 $\hat{y}_i$ 时才可使用。

\section{灵敏度分析}

灵敏度分析用于考察关键参数变化对模型输出的影响。若输出量为 $z$，参数为 $p$，可定义相对灵敏度：
\begin{equation}
  S_p = \frac{\Delta z / z}{\Delta p / p}.
  \label{eq:sensitivity}
\end{equation}
正式使用时需说明扰动范围、扰动方式和评价指标。

\section{结论}

结论应在完成数据处理、模型求解、检验和灵敏度分析后撰写。当前模板不包含真实数据和实验结果，因此不填写定量结论。

\section{模型优缺点}

\subsection{优点}
待根据模型结构、验证结果和实际表现填写。

\subsection{不足}
待根据模型假设、数据质量、误差来源和适用条件填写。


\section*{参考文献}
\addcontentsline{toc}{section}{参考文献}
当前模板没有真实引用，因此不列出参考文献。正式写作时，请从 Zotero 导出 BibTeX 到 \texttt{references.bib}，只引用已经实际阅读并用于论文的资料。

\appendix
\section{附录}
附录用于放置补充推导、伪代码、额外表格和代码说明。完整代码应保存在项目根目录的 \texttt{code/} 文件夹中。

\end{document}
